\documentclass[11pt]{article}

%------------------------------------------------
% Packages
%------------------------------------------------
\usepackage[a4paper,margin=2.5cm]{geometry}
\usepackage{amsmath,amssymb,amsthm}
\usepackage{mathtools}
\usepackage{enumitem}
\usepackage{hyperref}

\hypersetup{
    colorlinks=true,
    linkcolor=blue,
    citecolor=blue,
    urlcolor=blue
}

% Keep the bibliography compact enough for the one-page workshop format.
\let\oldthebibliography\thebibliography
\renewcommand{\thebibliography}[1]{%
  \small
  \oldthebibliography{#1}%
  \setlength{\itemsep}{0pt}%
  \setlength{\parskip}{0pt}%
}

%------------------------------------------------
% Title Information
%------------------------------------------------

\title{\bfseries An Invertible Generative Model for Forward and Inverse Problems}

\author{
Tristan van Leeuwen\thanks{Centrum Wiskunde \& Informatica, Amsterdam, and
Utrecht University, Utrecht, The Netherlands,
\texttt{T.van.Leeuwen@cwi.nl}}
\and
Christoph Brune\thanks{University of Twente, Enschede, The Netherlands}
\and
Marcello Carioni\footnotemark[2]
\and
Lasse Veenstra\thanks{Delft University of Technology, The Netherlands}
}

\date{}

%------------------------------------------------

\begin{document}

\maketitle

\vspace{-1.5em}

\noindent
\textbf{Keywords:}
Bayesian inverse problems, normalizing flows, triangular transport,
amortized inference

\vspace{-0.5em}

%%%%%%%%%%%%%%%%%%%%%%%%%%%%%%%%%%%%%%%%%%%%%%%%%%
\section*{Abstract}
\vspace{-0.5em}

Many inverse problems require both a forward simulator, which generates data
for given parameters, and an inference procedure, which characterizes the
parameters compatible with observed data. In a Bayesian formulation these
tasks amount to sampling from the likelihood and the posterior, respectively. 
Both tasks can benefit from recent advances in generative models. Existing amortized approaches typically learn separate conditional generative models for the two directions
\cite{Radev2023}, or focus primarily on the inverse direction
\cite{Ardizzone2019}.

We introduce a single invertible generative map that performs both tasks.
Starting from the classical use of triangular transport maps for conditional
sampling \cite{Marzouk2017}, we combine a lower-triangular transport, adapted
to the likelihood, with an upper-triangular transport, adapted to the
posterior \cite{vanLeeuwen2026}. The resulting bijection acts on the joint parameter--data space:
evaluated forward, it generates samples from the likelihood for a prescribed
parameter; evaluated in reverse, it generates posterior samples for prescribed
data. We investigate the (non-)uniqueness of this construction, and proposal a variational formulation that allows us to learn the map from samples.

In linear Gaussian problems the construction is explicit and,
unlike the individual triangular transports, can remain well conditioned as
the observational noise vanishes. Nonlinear examples and an image-inpainting
experiment illustrate simulation and uncertainty-aware inversion with the
same model. The framework thus offers a unified route to amortized forward
simulation and Bayesian inference, with potential extensions to experimental
design.

%%%%%%%%%%%%%%%%%%%%%%%%%%%%%%%%%%%%%%%%%%%%%%%%%%
\vspace{-0.5em}

\begin{thebibliography}{9}

\bibitem{Ardizzone2019}
L.~Ardizzone, J.~Kruse, S.~Wirkert, et al.
\newblock Analyzing inverse problems with invertible neural networks.
\newblock In {\em International Conference on Learning Representations}, 2019.

\bibitem{Marzouk2017}
Y.~Marzouk, T.~Moselhy, M.~Parno, and A.~Spantini.
\newblock Sampling via measure transport: An introduction.
\newblock In {\em Handbook of Uncertainty Quantification}, pp.~785--825,
Springer, 2017.

\bibitem{Radev2023}
S.~T. Radev, M.~Schmitt, V.~Pratz, et al.
\newblock JANA: Jointly amortized neural approximation of complex Bayesian
models.
\newblock In {\em Proceedings of UAI}, PMLR 216:1695--1706, 2023.

\bibitem{vanLeeuwen2026}
T.~van Leeuwen.
\newblock The reversible simulator: A data-driven approach for solving forward
and inverse problems.
\newblock In {\em Applications of Mathematics in Sciences, Engineering, and
Economics}, pp.~243--260, Springer, 2026.

\end{thebibliography}

\end{document}
